
\documentclass{article}
\usepackage{colortbl}
\usepackage{makecell}
\usepackage{multirow}
\usepackage{supertabular}

\begin{document}

\newcounter{utterance}

\centering \large Interaction Transcript for game `eqbench', experiment `eq\_bench', episode 82 with reflexion{-}llama{-}3.1{-}8b{-}instruct.
\vspace{24pt}

{ \footnotesize  \setcounter{utterance}{1}
\setlength{\tabcolsep}{0pt}
\begin{supertabular}{c@{$\;$}|p{.15\linewidth}@{}p{.15\linewidth}p{.15\linewidth}p{.15\linewidth}p{.15\linewidth}p{.15\linewidth}}
    \# & \multicolumn{2}{c}{Player} && \multicolumn{2}{c}{Game Master} \\
    \hline

    \theutterance \stepcounter{utterance}  
    & & & \multicolumn{4}{p{0.6\linewidth}}{
        \cellcolor[rgb]{0.9,0.9,0.9}{
            \makecell[{{p{\linewidth}}}]{
                \texttt{\tiny{[P1$\langle$GM]}}
                \texttt{You have attempted this same task before and did not succeed. These are the plans you wrote after those attempts.} \\
\texttt{Trial \#1: New plan: } \\
\\ 
\texttt{In my previous attempt, I scored the emotions based on my interpretation of the dialogue, but I didn't consider the context and the tone of the conversation. I should have taken into account that Charlie is being defensive and dismissive towards Mari, which suggests that he might be feeling indignant. However, I also scored him high on humiliation, which seems contradictory. } \\
\\ 
\texttt{Next time, I will pay closer attention to the tone and the language used in the dialogue. I will consider the power dynamics between the characters and the context of the conversation. For example, I might score Charlie higher on indignant and lower on humiliation, as he is more likely to feel defensive and dismissive rather than humiliated. I will also consider the fact that Mari is not being aggressive or condescending, which might suggest that Charlie is not feeling humiliated. } \\
\\ 
\texttt{I will also consider the fact that the dialogue is a conversation between two characters, and that the emotions expressed by one character might be influenced by the other character's words and tone. For example, Mari's comment "Seems like experience isn't everything" might be a subtle way of saying that Charlie's experience is not as valuable as he thinks it is, which could contribute to his indignant feeling. } \\
\\ 
\texttt{Overall, my new plan is to be more nuanced and context{-}dependent in my scoring, and to consider the tone, language, and power dynamics of the conversation.} \\
\\ 
\texttt{Trial \#2: New plan: } \\
\\ 
\texttt{In my previous attempt, I focused on the literal meaning of the dialogue, but I didn't consider the underlying dynamics and power struggles between the characters. I should have taken into account that Charlie is a seasoned police officer, while Mari is a younger and less experienced officer. This power imbalance could contribute to Charlie's indignant feeling, as he might feel that Mari is questioning his authority and experience.} \\
\\ 
\texttt{Next time, I will pay closer attention to the power dynamics and the tone of the conversation. I will consider the fact that Charlie is being defensive and dismissive towards Mari, which suggests that he might be feeling indignant. I will also consider the fact that Mari's comment "Seems like experience isn't everything" is a subtle way of saying that Charlie's experience is not as valuable as he thinks it is, which could contribute to his indignant feeling.} \\
\\ 
\texttt{To improve my scoring, I will ask myself the following questions:} \\
\\ 
\texttt{* What is the tone of the conversation? Is it confrontational, defensive, or dismissive?} \\
\texttt{* What are the power dynamics between the characters? Is there a power imbalance that could contribute to one character's emotions?} \\
\texttt{* How do the characters' words and tone influence each other's emotions?} \\
\texttt{* What are the underlying themes and messages in the dialogue? Are there any subtle hints or suggestions that could contribute to one character's emotions?} \\
\\ 
\texttt{By considering these factors, I hope to improve my scoring and provide a more nuanced and accurate interpretation of the dialogue.} \\
\\ 
\texttt{Trial \#3: New plan: } \\
\\ 
\texttt{In my previous attempts, I focused on the literal meaning of the dialogue and the emotions expressed by the characters, but I didn't consider the underlying dynamics and power struggles between the characters. I should have taken into account that Charlie is a seasoned police officer, while Mari is a younger and less experienced officer. This power imbalance could contribute to Charlie's indignant feeling, as he might feel that Mari is questioning his authority and experience.} \\
\\ 
\texttt{Next time, I will pay closer attention to the power dynamics and the tone of the conversation. I will consider the fact that Charlie is being defensive and dismissive towards Mari, which suggests that he might be feeling indignant. I will also consider the fact that Mari's comment "Seems like experience isn't everything" is a subtle way of saying that Charlie's experience is not as valuable as he thinks it is, which could contribute to his indignant feeling.} \\
\\ 
\texttt{To improve my scoring, I will ask myself the following questions:} \\
\\ 
\texttt{* What is the tone of the conversation? Is it confrontational, defensive, or dismissive?} \\
\texttt{* What are the power dynamics between the characters? Is there a power imbalance that could contribute to one character's emotions?} \\
\texttt{* How do the characters' words and tone influence each other's emotions?} \\
\texttt{* What are the underlying themes and messages in the dialogue? Are there any subtle hints or suggestions that could contribute to one character's emotions?} \\
\\ 
\texttt{By considering these factors, I hope to improve my scoring and provide a more nuanced and accurate interpretation of the dialogue.} \\
\\ 
\\ 
\\ 
\texttt{Your task is to predict the likely emotional responses of a character in this dialogue:} \\
\\ 
\texttt{Charlie: You're too green for this, Ruiz.} \\
\texttt{Mari: And you're too washed up, Redburn.} \\
\texttt{Charlie: I've been on this beat longer than you've been alive, kid.} \\
\texttt{Mari: And yet, here we are, stuck on the same case. Seems like experience isn't everything.} \\
\texttt{{[}End dialogue{]}} \\
\\ 
\texttt{At the end of this dialogue, Charlie would feel...} \\
\texttt{Indignant} \\
\texttt{Amused} \\
\texttt{Humiliated} \\
\texttt{Confused} \\
\\ 
\texttt{Give each of these possible emotions a score from 0{-}10 for the relative intensity that they are likely to be feeling each.} \\
\\ 
\texttt{You must output in the following format, including headings (of course, you should give your own scores), with no additional commentary:} \\
\\ 
\texttt{Indignant: <score>} \\
\texttt{Amused: <score>} \\
\texttt{Humiliated: <score>} \\
\texttt{Confused: <score>} \\
\\ 
\\ 
\texttt{{[}End of answer{]}} \\
\\ 
\texttt{Remember: zero is a valid score, meaning they are likely not feeling that emotion. You must score at least one emotion > 0.} \\
\\ 
\texttt{Your answer:} \\
            }
        }
    }
     \\ \\

    \theutterance \stepcounter{utterance}  
    & \multicolumn{4}{p{0.6\linewidth}}{
        \cellcolor[rgb]{0.561,0.737,0.561}{
            \makecell[{{p{\linewidth}}}]{
                \texttt{\tiny{[P1$\rangle$GM]}}
                \texttt{Indignant: 8} \\
\texttt{Amused: 0} \\
\texttt{Humiliated: 2} \\
\texttt{Confused: 0} \\
            }
        }
    }
    & & \\ \\

    \theutterance \stepcounter{utterance}  
    & & & \multicolumn{2}{p{0.3\linewidth}}{
        \cellcolor[rgb]{0.9,0.9,0.9}{
            \makecell[{{p{\linewidth}}}]{
                \texttt{\tiny{[GM$|$GM]}}
                \texttt{Indignant: 8} \\
\texttt{Amused: 0} \\
\texttt{Humiliated: 2} \\
\texttt{Confused: 0} \\
            }
        }
    }
    & & \\ \\

    \theutterance \stepcounter{utterance}  
    & & & \multicolumn{2}{p{0.3\linewidth}}{
        \cellcolor[rgb]{0.9,0.9,0.9}{
            \makecell[{{p{\linewidth}}}]{
                \texttt{\tiny{[GM$|$GM]}}
                \texttt{Indignant: 8} \\
\texttt{Amused: 0} \\
\texttt{Humiliated: 5} \\
\texttt{Confused: 0} \\
            }
        }
    }
    & & \\ \\

    \theutterance \stepcounter{utterance}  
    & & & \multicolumn{2}{p{0.3\linewidth}}{
        \cellcolor[rgb]{0.9,0.9,0.9}{
            \makecell[{{p{\linewidth}}}]{
                \texttt{\tiny{[GM$|$GM]}}
                \texttt{game\_result = LOSE} \\
            }
        }
    }
    & & \\ \\

    \theutterance \stepcounter{utterance}  
    & & \multicolumn{4}{p{0.6\linewidth}}{
        \cellcolor[rgb]{0.95,0.95,0.95}{
            \makecell[{{p{\linewidth}}}]{
                \texttt{\tiny{[P1$\leftrightarrow$P]}}
                \texttt{Post{-}episode Reflexion (not sent to GM):} \\
\texttt{New plan: } \\
\\ 
\texttt{In my previous attempts, I focused on the literal meaning of the dialogue and the emotions expressed by the characters, but I didn't consider the underlying dynamics and power struggles between the characters. I should have taken into account that Charlie is a seasoned police officer, while Mari is a younger and less experienced officer. This power imbalance could contribute to Charlie's indignant feeling, as he might feel that Mari is questioning his authority and experience.} \\
\\ 
\texttt{However, I also realize that I should have considered the fact that Mari's comment "Seems like experience isn't everything" is not necessarily a personal attack on Charlie's experience, but rather a general statement about the importance of experience in a particular situation. This could suggest that Charlie's indignant feeling is not as strong as I initially thought.} \\
\\ 
\texttt{To improve my scoring, I will ask myself the following questions:} \\
\\ 
\texttt{* What is the tone of the conversation? Is it confrontational, defensive, or dismissive?} \\
\texttt{* What are the power dynamics between the characters? Is there a power imbalance that could contribute to one character's emotions?} \\
\texttt{* How do the characters' words and tone influence each other's emotions?} \\
\texttt{* What are the underlying themes and messages in the dialogue? Are there any subtle hints or suggestions that could contribute to one character's emotions?} \\
\texttt{* How does Mari's comment "Seems like experience isn't everything" relate to the conversation and Charlie's emotions?} \\
\\ 
\texttt{By considering these factors, I hope to improve my scoring and provide a more nuanced and accurate interpretation of the dialogue.} \\
\\ 
\texttt{Indignant: 9} \\
\texttt{Amused: 0} \\
\texttt{Humiliated: 1} \\
\texttt{Confused: 0} \\
            }
        }
    }
    & \\ \\

\end{supertabular}
}

\end{document}
